%% Section 6: The Shuffle Algebra Substrate
%% To be \input{} into the main document.

\section{The Shuffle Algebra Substrate}
\label{sec:shuffle}

The three-level hierarchy of Section~\ref{sec:hierarchy} organises the
space of LR models vertically (by integrability level) and
horizontally (by deformation type).  But the hierarchy itself sits
inside a larger algebraic universe.  In this section we sketch the
contours of that universe: the \emph{shuffle algebra} framework, which
provides a single algebraic substrate from which all the integrable
vertex models of this paper can be derived.

\subsection{Feigin-Odesskii shuffle algebras}
\label{sec:shuffle-def}

Let $V$ be a finite-dimensional vector space and let
$\zeta(z) \in \mathbb{C}(z)$ be a rational function, called the
\emph{shuffle kernel}.  The \emph{Feigin-Odesskii shuffle
algebra}~\cite{feigin-odesskii} is defined on the direct sum
\[
  \mathcal{S}_\zeta \;=\; \bigoplus_{n \geq 0} \mathbb{C}(z_1, \ldots, z_n),
\]
with multiplication given by the \emph{shuffle product}: for
$f \in \mathbb{C}(z_1, \ldots, z_m)$ and
$g \in \mathbb{C}(z_1, \ldots, z_n)$,
\begin{equation}\label{eq:shuffle-product}
  (f \star g)(z_1, \ldots, z_{m+n})
  \;=\;
  \mathrm{Sym}\!\left[
    f(z_1, \ldots, z_m) \cdot g(z_{m+1}, \ldots, z_{m+n})
    \cdot \prod_{i=1}^{m} \prod_{j=m+1}^{m+n}
    \zeta\!\left(\frac{z_i}{z_j}\right)
  \right],
\end{equation}
where $\mathrm{Sym}$ denotes symmetrisation over all
$(m,n)$-shuffles of the variables $z_1, \ldots, z_{m+n}$.  The
shuffle kernel $\zeta$ controls the algebra: different choices of
$\zeta$ produce genuinely different algebraic structures.  At the
extreme $\zeta \equiv 1$, the shuffle product reduces to ordinary
symmetrisation and $\mathcal{S}_\zeta$ is the algebra of symmetric
rational functions.  Nontrivial kernels introduce interaction between
the two factors, and it is precisely this interaction that encodes the
Boltzmann weights of integrable lattice models.

\subsection{From shuffle algebras to vertex models}

The connection to our story was established by Garbali and
Gunna~\cite{garbali-gunna-2024}, who proved that the vertex models
computing symmetric function structure constants---including the
Schur, Grothendieck, and Hall-Littlewood models of
Section~\ref{sec:lattice}---can be realised as \emph{representations}
of appropriate shuffle algebras.  More precisely, the row operators of
the lattice model (the transfer matrices $T(u)$ of
Section~\ref{sec:lattice-transfer}) intertwine the shuffle product
with the composition of vertex weights: the algebraic identity
$T(u) \star T(v) = T(u) \circ T(v)$ (suitably interpreted) encodes
both the Yang-Baxter equation and the commutativity of transfer
matrices in a single stroke.

Garbali and Negut~\cite{garbali-negut-2025} extended this picture by
proving that the shuffle algebra $\mathcal{S}_\zeta$, for the kernel
$\zeta(z) = (z - 1)/(z - q)$ relevant to the Schur/Hall-Littlewood
family, is isomorphic to the positive part $U_q^+(\widehat{\mathfrak{gl}}_1)$
of the quantum toroidal algebra (equivalently, the positive half of
the Yangian $Y(\widehat{\mathfrak{gl}}_1)$ at the rational level).
This isomorphism is not merely formal: it identifies the shuffle
generators with specific combinations of Drinfeld currents, and under
this identification, the vertex model partition functions become
matrix elements of intertwining operators between Fock space modules.
The consequence is stark: \emph{all} the integrable lattice models
from Section~\ref{sec:lattice}---Schur, Grothendieck,
Hall-Littlewood, $q$-Whittaker---are manifestations of a single
algebraic structure, viewed through different representations.

\subsection{Connection to the Hecke algebra}

The shuffle algebra also makes contact with the Hecke transition
algebra of Section~\ref{sec:hierarchy}.  Kalmykov~\cite{kalmykov}
established explicit isomorphisms connecting three objects at the
rational level:
\begin{enumerate}
\item the Feigin-Odesskii shuffle algebra $\mathcal{S}_\zeta$;
\item the \emph{degenerate affine Hecke algebra} (dAHA)
  $H_n^{\mathrm{aff}}$, which extends the finite Hecke algebra $H_q(S_n)$
  by polynomial generators $x_1, \ldots, x_n$ satisfying the
  Bernstein-Lusztig relations; and
\item the Yangian $Y(\mathfrak{gl}_n)$.
\end{enumerate}
Since the finite Hecke algebra $H_q(S_n)$ embeds in the dAHA as the
subalgebra generated by the simple transpositions (without the
polynomial generators), the Hecke transition algebra from
Section~\ref{sec:hierarchy} sits inside the shuffle substrate as a
natural subalgebra.  The three-level hierarchy
$H_q \to H_0 \to \mathbb{C}[S_n]$ is a shadow of a deeper hierarchy
within the Yangian, where the specialisation $q \to 0$ corresponds to
a contraction of the quantum group to its classical limit.

\subsection{The emerging picture}

The shuffle algebra framework suggests a four-level depth hierarchy
for understanding why the LR coefficients admit so many combinatorial
models:

\begin{enumerate}
\item \textbf{Combinatorial bijections.}  At the surface, one has
  explicit maps between tableaux, puzzles, pipe dreams, and other
  models.  These bijections are typically proved by induction or
  sign-reversing involutions, and their existence appears miraculous.

\item \textbf{Yang-Baxter groupoid.}  One level down, the bijections
  are morphisms in the Yang-Baxter groupoid of Bump and
  Naprienko~\cite{bump-naprienko-2025}
  (Section~\ref{sec:groupoid-bn}).  Their existence follows from the
  Yang-Baxter equation: each R-matrix move preserves the partition
  function, so composing moves gives a canonical bijection.  The
  miracle is demystified but not fully explained---one still needs to
  know \emph{why} the R-matrix satisfies the YBE.

\item \textbf{Shuffle algebra.}  At the next level, the R-matrices
  are intertwining operators for representations of the shuffle
  algebra $\mathcal{S}_\zeta$.  The YBE is a consequence of the
  associativity of the shuffle product, and the multiplicity of
  lattice models reflects the multiplicity of representations.
  Different representations give different vertex weights, but the
  partition functions agree because the underlying algebra is the
  same.

\item \textbf{Geometric Satake.}  At the deepest level (which we do
  not develop here), the shuffle algebra is the cohomology of a
  moduli space of instantons, and the LR coefficients are intersection
  numbers on the affine Grassmannian.  The geometric Satake
  correspondence~\cite{mirkovic-vilonen-2007} provides the ultimate
  explanation: the structure constants of the representation ring are
  geometric invariants, and their positivity is a consequence of the
  decomposition theorem in perverse sheaves.
\end{enumerate}

\noindent
Each level subsumes the ones above it, offering a deeper explanation
for the same numerical coincidences.  The categorical phase diagram of
Section~\ref{sec:hierarchy} lives at levels~2 and~3: the vertical
axis (integrability level) is governed by the Hecke quotients of the
shuffle algebra, while the horizontal axis (deformation type) is
governed by the choice of shuffle kernel~$\zeta$.  The full picture---a
phase diagram sitting inside the representation theory of a shuffle
algebra, itself a shadow of geometric Satake---remains to be completed.
We turn to the open questions in the next section.
